\section{Introduction}

\subsection{The option contract}

I am interested in pricing fixed-strike, arithmetic-average Asian call options.  The
payoff at expiry $T$ is
\begin{equation}
  C(T) = \max\!\left(A(0,T) - K,\, 0\right), \qquad
  A(0,T) = \frac{1}{N}\sum_{j=1}^{N} S_{t_j},
\end{equation}
where $S_{t_j}$ is the price of the underlying at time $t_j = jT/N$ and $K$ is the
strike.  Path-dependence through the average $A$ prevents the Black--Scholes machinery
from yielding an exact closed form, making approximations or simulation necessary.
\citet{levy1992pricing} proposed a moment-matching approximation: matching the first
two moments of the discrete arithmetic average to a lognormal yields a tractable
formula that is accurate under constant-volatility Black--Scholes dynamics, where $S_t$
is lognormal.

The situation is more difficult under Rough Fractional Stochastic Volatility (RFSV)
\citep{gatheral2018volatility}, in which $\log\sigma_t = \nu W_t^H$, where $W_t^H$ is
a fractional Brownian motion with Hurst exponent $H \approx 0.10$.  Under RFSV, $S_t$
is no longer lognormal, and the two-moment Lévy approximation cannot capture the heavy
skew and fat tails introduced by a large vol-of-vol $\nu$.  Computing even the second
joint moment $\mathbb{E}[S_{t_i}S_{t_j}]$ analytically is intractable: it requires the
expectation of the integrated volatility path, which depends on the full continuous
history of $W^H$ up to each time point.

Monte Carlo sidesteps the distributional intractability by estimating the expectation
directly from simulated paths.  By the law of large numbers,
\begin{equation}
  \frac{1}{M}\sum_{m=1}^{M} e^{-rT}\max\!\left(A^{(m)} - K,\,0\right)
  \xrightarrow{M \to \infty}
  \mathbb{E}^{\mathbb{Q}}\!\left[e^{-rT}\max(A - K,\, 0)\right],
\end{equation}
where $\mathbb{Q}$ is the risk-neutral measure and $A^{(m)}$ is the arithmetic average
along the $m$-th simulated path.  No closed-form density or moment-matching is required;
we need only the ability to simulate correlated log-volatility paths.

\subsection{Why rough volatility}

\citet{gatheral2018volatility} provide compelling empirical evidence that equity
log-realized-variance behaves as fractional Brownian motion with $H \approx 0.10$.
Estimating $H$ by ordinary least squares on log-realized-variance increments across
five major indices (S\&P 500, NASDAQ, DAX, FTSE 100, Nikkei 225) and time scales
ranging from one day to one month, they find a consistent Hurst exponent well below
$\nicefrac{1}{2}$, with $R^2 > 0.90$ in each case.  I take $H = 0.10$ and their
vol-of-vol estimate from high-frequency realized variance data \citep{oxfordman},
$\nu \approx 0.3$.  One subtlety concerns units: \citet{gatheral2018volatility}
measure time in days, while my engine measures it in years ($T = 1$).  Because the
log-vol increment variance scales as $\nu^2\Delta^{2H}$, the same fit in year units is
$\nu = 0.3 \cdot 252^{0.1} \approx 0.52$, which is the value I use.  My own variogram fit to
non-overlapping 5-day realized variance from \texttt{yfinance}, with time in years,
gives $H \approx 0.12$ ($R^2 \approx 0.98$ across a handful of lags) and
$\nu \approx 0.87$, the latter inflated by measurement noise in weekly
squared-return variance.

Standard Brownian motion corresponds to $H = \nicefrac{1}{2}$, with uncorrelated increments. 
At $H = 0.10$, volatility paths are far rougher than BM: increments
are negatively correlated at all lags, the path is less regular than classical
stochastic volatility models, and the realized variance process mean-reverts on
short timescales. The short-dated implied volatility smile
steepens beyond what classical stochastic volatility models expect, and
volatility clustering manifests on timescales an order of magnitude shorter than, e.g., Heston.

The RFSV model I studied is \citep{gatheral2018volatility}
\begin{equation}
  \log\sigma_t = \log\sigma_0 + \nu W_t^H, \qquad
  dS_t = r S_t\,dt + \sigma_t S_t\, dW_t^{\mathbb{Q}},
\end{equation}
where $W_t^H$ is the fractional Brownian motion of \citet{mandelbrot1968fractional}
with covariance $\mathbb{E}[W_t^H W_s^H] = \tfrac{1}{2}(t^{2H}+s^{2H}-|t-s|^{2H})$.
I set $r = 0$ throughout for simplicity (consistent with \texttt{params.hpp}); the
drift term drops out of all pricing formulas but is retained in the SDE for
completeness.
I use $H = 0.10$, $\nu = 0.52$, and base volatility $\sigma_0 = 0.2$ throughout, and
also explore the sensitivity of option prices to $H$ and $\nu$ in
Section~\ref{sec:sensitivity}.  The base level is not a cosmetic choice.  An earlier
version of this study used $\sigma_0 = 1$ (100\% volatility); at $\nu = 0.52$ that
payoff is so heavy-tailed that a single path in $10^6$ carried 34\% of the sample
variance, and Monte Carlo error bars were meaningless.  At $\sigma_0 = 0.2$ the largest
path carries $0.11\%$, and the payoff standard deviation is stable across batches.
Strictly, the payoff variance is infinite for any lognormal volatility model, because
$\mathbb{E}[S_t^2]$ contains $\mathbb{E}[\exp(\Delta t\, e^{2\nu W})]$; in practice the
divergence sits so far in the tail that it is never sampled.

\subsection{The computational bottleneck}
\label{sec:bottleneck}

Naïve Monte Carlo for RFSV requires $O(N^3)$ setup work.  Because fBM is
non-Markovian, there is no forward recursion: the value of $W_t^H$ depends on the
entire history of the driving noise, not just its most recent value.  Generating each
sample path therefore requires drawing from the full $N$-dimensional Gaussian
$\mathcal{N}(0, C)$, where $C$ is a dense $N \times N$ positive-definite matrix with
$C_{ij} = \tfrac{1}{2}(t_i^{2H}+t_j^{2H}-|t_i-t_j|^{2H})$.  The standard approach
factorizes $C = LL^\top$ via Cholesky decomposition once ($O(N^3/3)$ flops), then
generates each path by computing $\nu Lz$ for $z \sim \mathcal{N}(0,I)$ ($O(N^2)$
flops per path, a lower-triangular matrix-vector multiply).

At $N = 252$ trading days and $M = 10{,}000$ paths the costs are concrete.  Cholesky
decomposition requires exactly $N^3/3 \approx 5.3$ million floating-point operations
(the leading constant $\nicefrac{1}{3}$ compares favourably to $\nicefrac{2}{3}$ for
LU); this is a one-time cost.  The per-path triangular multiply costs $O(N^2) \approx 63{,}000$ flops, so
the total Monte Carlo loop costs $O(MN^2) \approx 630$ million flops, more than
$100\times$ the factorization.  MC pricing error falls as $O(M^{-1/2})$: with
$\sigma_\mathrm{payoff} \approx 9.5$ (Section~\ref{sec:convergence}), $M = 10{,}000$
paths give a standard error of $\approx 0.095$ price units, about $1.8\%$ of the
at-the-money price of $\approx 5.3$, and halving it requires four times as many paths.

Trading desks building sensitivity surfaces run parameter sweeps over $H$, $\nu$, and
$K$, easily multiplying $M$ by factors of $10$--$100$.  Intraday risk monitoring at
$N \approx 19{,}656$ (five-minute intervals over 252 trading days of 6.5 hours each:
$252 \times 78 = 19{,}656$) increases the per-path cost by $\sim 6{,}000\times$
relative to daily steps.

\subsection{Approach and scope}
\label{sec:scope}

I implement and benchmark three algorithms that reduce the cost of fBM path simulation
by exploiting structure in the covariance matrix $C$:

\begin{enumerate}
  \item \textbf{Dense Cholesky} ($O(N^2)$/path, exact): exploits the
        positive-definiteness of $C$ to factor it exactly.  By ``exact simulation''
        I mean that the method is free of time-discretization bias: it samples the correct
        fBM distribution for any finite $N$, unlike Euler--Maruyama approximations.
        Serves as the correctness baseline against which the other methods are validated.

  \item \textbf{Circulant embedding \& FFT} ($O(N\log N)$/path): exploits the
        fact that the increments of fBM (fractional Gaussian noise, fGn) are
        stationary. Their covariance matrix is Toeplitz, which can be embedded
        into a circulant matrix, which can be diagonalized by the DFT
        \citep{davies1987tests, wood1994simulation}.  For $H \leq \nicefrac{1}{2}$
        the embedding is guaranteed positive \citep{craigmile2003simulating}, so the
        method is exact on the grid.

  \item \textbf{Global low-rank rSVD} ($O(Nk)$/path): exploits the low-rank
        structure of smooth off-diagonal blocks of $C$: the off-diagonal blocks admit low-rank
        approximations.  A randomized SVD \citep{halko2011finding} compresses $C
        \approx U\,\mathrm{diag}(S)\,U^\top$ at rank $k$, reducing per-path cost from
        $O(N^2)$ to $O(Nk)$.  In practice the implementation is a global
        approximation of the entire $N\times N$ matrix (including the singular
        near-diagonal region) rather than of the smooth off-diagonal blocks alone.
        This is blunter, and must allocate rank budget to the diagonal singularity,
        which is why the sampler still discards $13\%$ of the path variance at
        $k = 128$ for $H = 0.10$, and underprices by $\approx 4\%$
        (§\ref{sec:errorrank}).
\end{enumerate}

Section~\ref{sec:foundations} develops the mathematical properties of $C$ that each
method exploits.  Section~\ref{sec:algorithms} describes each algorithm and its
complexity.  Sections~\ref{sec:benchmarks} and~\ref{sec:validation} present empirical
timing, memory, accuracy, and model validation results.

The project does not address variance reduction (quasi-Monte Carlo, control variates),
model extensions beyond fBM-RFSV (rough Bergomi, multifactor), or GPU parallelism,
though Section~\ref{sec:futurework} situates each of these as a natural next step.
